This chapter will cover the follow-up research to be done, using the results provided above. Following the results presented above, the next steps in the analysis will be as follows:

\begin{description}
\item [ Evaluation of systematics ] The results given above for the accuracies have only statistical uncertainties applied to them. Therefore, in order to have good control over the results, method-dependent systematic uncertainties and closure tests evaluated.

\item [Extracting physics from the ML model] One of the reasons for creating the ML model as built was to extract new observables and identify differences between the samples provided. As mentioned in section \ref{sec:ML}, the inclusion of the ResNet models allows us to look into the decision making process of the model, and analyze the prominent filters in a single event. Given the results of the model, it is clear there is a latent space in which it is possible to distinguish between different photon induced processes , resulting in di-jets. Therefore, identifying these observables could lead to better understanding of such processes.

\item [Processing Run 3 data in order to increase statistics for future measurements] As mentioned, The results previously used only the 2018 Pb + Pb data. In addition, the reduction in number of events due to event selection, as seen in section \ref{sec:event_selection}, means that there is a lack of statistics to achieve precision. Thus, one of the important tasks is to increase the number of events available by going through and processing Run 3 UPC di-jet events, an effort which will include participating in the jet calibration process in ATLAS for all the Run 3 datasets, as well as other tasks that may arise in getting the data to the position in which the same analysis can be run on.

\item [Validation of the $\gamma \mathbb{P}$ MC sample] While we use the $\gamma \mathbb{P}$ MC in this analysis, there are some open questions regarding its handling of neutrals, etc. In addition, there is currently no other MC to compare against, which has been processed through the entire ATLAS analysis chain. Therefore another task is to conduct a detailed validation of the MC.
\end{description}


In addition to all of these tasks, There is also work being done on the ZDC (The sub detector mentioned in section \ref{sec:atlas}). As part of the ZDC team and an expert on ZDC TDAQ infrastructure, There is additional work that will carry on with respect to the last runs of Run 3 and preparation for Run 4. In contrast to the amount of information presented here, work on the ZDC takes considerable amount of time, but the result is better data taking that can be used for future measurements.
